\begin{tikzpicture}[
    scale=0.8,
    transform shape,
    every node/.style={
        draw,
        rounded corners,
        align=center,
        text width=2.0cm,
        minimum height=0.6cm,
        inner sep=2pt,
        font=\scriptsize
    },
    main/.style={fill=blue!40, font=\bfseries},
    secondary/.style={fill=blue!20},
    neutral/.style={fill=gray!15},
    arrow/.style={->, thick}
]

% ===== NIVEL 0 =====
\node[main, text width=5cm,font=\small] (main) at (0,0) {Procesamiento de señales en sistemas embebidos};

% ===== NIVEL 1 =====
\node[secondary] (rep)   at (-8,-2) {Representación de señales};
\node[secondary] (mod)   at (-2.5,-2)  {Modulación de señales};
\node[secondary] (demod) at (2.5,-2)  {Demodulación de señales};
\node[secondary] (emb)   at (8.5,-2)  {Sistemas embebidos};
     
% ===== NIVEL 2 =====
\node[neutral] (time) at (-9.5,-4) {Dominio del tiempo};
\node[neutral] (freq) at (-6.5,-4) {Dominio de la frecuencia};

\node[neutral] (modA) at (-3.7,-4) {Modulación analógica};
\node[neutral] (modD) at (-1.2,-4)  {Modulación digital};

\node[neutral] (env) at (1.2,-4) {Demodulación por envolvente};
\node[neutral] (coh) at (3.7,-4) {Demodulación coherente};

\node[neutral] (rt)   at (6,-4) {Ejecución en tiempo real};
\node[neutral] (hw)   at (8.5,-4) {Restricciones de hardware};
\node[neutral] (dsp)   at (11,-4) {Implementación de DSP};

% ===== NIVEL 3 =====
\node[neutral] (analog)  at (-9.5,-6.5) {Analógica};
\node[neutral] (digital) at (-6.5,-6.5) {Digital};

\node[neutral] (proc) at (7.3,-6) {Capacidad de procesamiento};
\node[neutral] (mem)  at (9.9,-6) {Memoria};

\node[secondary] (impl)  at (0,-6) {Implementación digital};


% ===== NIVEL 4 =====

\node[neutral] (sample) at (-8,-8.5) {Muestreo};
\node[neutral] (disc)   at (-3,-8) {Señales discretas};
\node[neutral] (block)  at (0,-8)  {Procesamiento por bloques};
\node[neutral] (filter) at (3,-8)  {Filtrado digital};

% ===== CONEXIONES ORTOGONALES =====

% Nivel 0
\draw[arrow] (main) -| (rep);
\draw[arrow] (main) |- (mod);
\draw[arrow] (main) |- (demod);
\draw[arrow] (main) -| (emb);

% Representación
\coordinate (split_rep) at ($(rep)+(0,-1)$);

\draw (rep) -- (split_rep);

\draw (split_rep) -- ($(time |- split_rep)$);
\draw (split_rep) -- ($(freq |- split_rep)$);

\draw[arrow] ($(time |- split_rep)$) -- (time);
\draw[arrow] ($(freq |- split_rep)$) -- (freq);

% ===== PRIMER NIVEL (time + freq → barra superior) =====

\coordinate (top_merge) at ($(time)!0.5!(freq)+(0,-1)$);

% Bajadas
\draw (time) -- ($(time |- top_merge)$);
\draw (freq) -- ($(freq |- top_merge)$);

% Barra horizontal
\draw ($(time |- top_merge)$) -- ($(freq |- top_merge)$);


% ===== SEGUNDO NIVEL (barra superior → barra inferior) =====

\coordinate (bottom_split) at ($(top_merge)+(0,-0.5)$);

\draw (top_merge) -- (bottom_split);

% Barra inferior
\draw (bottom_split) -- ($(analog |- bottom_split)$);
\draw (bottom_split) -- ($(digital |- bottom_split)$);


% ===== SALIDAS =====

\draw[arrow] ($(analog |- bottom_split)$) -- (analog);
\draw[arrow] ($(digital |- bottom_split)$) -- (digital);

% Punto de unión arriba de sample
\coordinate (merge_sample) at ($(sample)+(0,1)$);

% Bajadas desde analog y digital hacia la misma altura
\draw (analog) -- ($(analog |- merge_sample)$);
\draw (digital) -- ($(digital |- merge_sample)$);

% Barra horizontal
\draw ($(analog |- merge_sample)$) -- ($(digital |- merge_sample)$);

% Flecha final
\draw[arrow] (merge_sample) -- (sample);


% Modulación
% Punto debajo de mod
\coordinate (split_mod) at ($(mod)+(0,-1)$);

% Bajada sin flecha
\draw (mod) -- (split_mod);

% Barra horizontal
\draw (split_mod) -- ($(modA |- split_mod)$);
\draw (split_mod) -- ($(modD |- split_mod)$);

% Flechas hacia abajo
\draw[arrow] ($(modA |- split_mod)$) -- (modA);
\draw[arrow] ($(modD |- split_mod)$) -- (modD);

% Demodulación
% Punto intermedio
\coordinate (split_demod) at ($(demod)+(0,-1)$);

% Bajada sin flecha
\draw (demod) -- (split_demod);

% Barra horizontal
\draw (split_demod) -- ($(env |- split_demod)$);
\draw (split_demod) -- ($(coh |- split_demod)$);

% Flechas hacia abajo
\draw[arrow] ($(env |- split_demod)$) -- (env);
\draw[arrow] ($(coh |- split_demod)$) -- (coh);


% ===== BARRA DESDE MODULACIÓN Y DEMODULACIÓN HACIA IMPL =====

% Nivel de la barra (ajusta altura si quieres)
\coordinate (bar_impl) at ($(impl)+(0,1)$);

% Bajadas desde los 4 nodos
\draw (modA) -- ($(modA |- bar_impl)$);
\draw (modD) -- ($(modD |- bar_impl)$);
\draw (env)  -- ($(env  |- bar_impl)$);
\draw (coh)  -- ($(coh  |- bar_impl)$);

% Barra horizontal completa
\draw ($(modA |- bar_impl)$) -- ($(coh |- bar_impl)$);

% Flecha final al nodo
\draw[arrow] (bar_impl) -- (impl);


% Sistemas embebidos
% ===== SPLIT DESDE emb (3 salidas) =====

\coordinate (split_emb) at ($(emb)+(0,-1)$);

% Bajada
\draw (emb) -- (split_emb);

% Barra horizontal (3 ramas)
\draw (split_emb) -- ($(rt |- split_emb)$);
\draw (split_emb) -- ($(hw |- split_emb)$);
\draw (split_emb) -- ($(dsp |- split_emb)$);

% Flechas hacia nodos
\draw[arrow] ($(rt |- split_emb)$) -- (rt);
\draw[arrow] ($(hw |- split_emb)$) -- (hw);
\draw[arrow] ($(dsp |- split_emb)$) -- (dsp);




% ===== SPLIT DESDE hw =====

\coordinate (split_hw) at ($(hw)+(0,-1)$);

% Bajada
\draw (hw) -- (split_hw);

% Barra horizontal
\draw (split_hw) -- ($(proc |- split_hw)$);
\draw (split_hw) -- ($(mem |- split_hw)$);

% Flechas finales
\draw[arrow] ($(proc |- split_hw)$) -- (proc);
\draw[arrow] ($(mem |- split_hw)$) -- (mem);



% Punto debajo de impl (baja 1 unidad)
\coordinate (split) at ($(impl)+(0,-1)$);

% Línea principal
\draw (impl) -- (split);

% Barra horizontal
\draw (split) -- ($(disc |- split)$);
\draw (split) -- ($(filter |- split)$);

% Bajadas
\draw[arrow] ($(disc |- split)$) -- (disc);
\draw[arrow] (split) -- (block);
\draw[arrow] ($(filter |- split)$) -- (filter);

\end{tikzpicture}